\section*{Appendix}

The appendix is organized to keep reproducibility and diagnostic detail outside the main narrative:
\begin{enumerate}[leftmargin=*]
    \item Appendix \ref{app:sec:formal-details}: formal details and notation.
    \item Appendix \ref{app:sec:atm-details}: ATM-Bench data, prompts, evidence views, and evaluation details.
    \item Appendix \ref{app:sec:stale-details}: STALE data, primitive-role mapping, and evaluation details.
    \item Appendix \ref{app:sec:error-analysis}: qualitative error taxonomy.
    \item Appendix \ref{app:sec:additional-results}: additional tables and robustness checks.
\end{enumerate}

\clearpage

\setcounter{secnumdepth}{1}
\renewcommand{\thesection}{\Alph{section}}
\setcounter{section}{0}

\section{Formal Details}
\label{app:sec:formal-details}

This appendix gives the full proofs for the two formal claims in Section~\ref{sec:framework}.

\subsection{Proof of View Collapse Impossibility}
\label{app:proof-view-collapse}

Let $D$ be any decoder whose input is only the view $V_q(\mathcal A)$.
Suppose
\[
V_q(\mathcal A_1)=V_q(\mathcal A_2)
\]
and
\[
\mathcal Z_q(\mathcal A_1)\cap\mathcal Z_q(\mathcal A_2)=\emptyset .
\]
For a deterministic decoder, the identical view implies that the decoder outputs the same response $y$ in both evidence worlds.
If $y$ is evidence-valid for $\mathcal A_1$, then
\[
\eta(y)\in \mathcal Z_q(\mathcal A_1).
\]
If the same $y$ is evidence-valid for $\mathcal A_2$, then
\[
\eta(y)\in \mathcal Z_q(\mathcal A_2).
\]
Thus $\eta(y)\in \mathcal Z_q(\mathcal A_1)\cap\mathcal Z_q(\mathcal A_2)$, contradicting the disjointness assumption.
Therefore, no deterministic view-only decoder can be evidence-valid on both worlds.

For a randomized decoder, the same view induces the same output distribution $p(y\mid V_q)$ in both worlds.
To be evidence-valid with probability one on $\mathcal A_1$, its support must be contained in
\[
\{y:\eta(y)\in\mathcal Z_q(\mathcal A_1)\}.
\]
To be evidence-valid with probability one on $\mathcal A_2$, the same support must be contained in
\[
\{y:\eta(y)\in\mathcal Z_q(\mathcal A_2)\}.
\]
Because $\eta$ is single-valued where defined and the semantic answer sets are disjoint, these two valid-output sets are disjoint.
No nonempty probability distribution can have support contained in both.
Hence a randomized view-only decoder also cannot be evidence-valid on both worlds with probability one.

\subsection{Proof of Dynamic Propagation of Binding Error}
\label{app:proof-dynamic-binding}

We first bound the one-step binding perturbation.
Since all values in $S_t$ equal $v_t^\star$,
\[
\tilde v_t
=
\sum_{i\in S_t}\alpha_{ti}v_i
+
\sum_{i\notin S_t}\alpha_{ti}v_i
=
B_t v_t^\star
+
\sum_{i\notin S_t}\alpha_{ti}v_i .
\]
Using $\sum_i\alpha_{ti}=1$, we have
\[
\tilde v_t-v_t^\star
=
\sum_{i\notin S_t}\alpha_{ti}(v_i-v_t^\star).
\]
Taking norms and applying the triangle inequality gives
\[
\lVert\tilde v_t-v_t^\star\rVert_2
\le
\sum_{i\notin S_t}\alpha_{ti}\lVert v_i-v_t^\star\rVert_2 .
\]
By the definition of $D_t$, $\lVert v_i-v_t^\star\rVert_2\le D_t$ for all $i$.
Hence
\[
\lVert\tilde v_t-v_t^\star\rVert_2
\le
\sum_{i\notin S_t}\alpha_{ti}D_t
=
(1-B_t)D_t .
\]

Now consider the evaluation-state error:
\[
e_{t+1}
=
\lVert h_{t+1}-h_{t+1}^\star\rVert_2
=
\lVert T_t(h_t,\tilde v_t)-T_t(h_t^\star,v_t^\star)\rVert_2 .
\]
By the local Lipschitz condition,
\[
e_{t+1}
\le
\rho_t\lVert h_t-h_t^\star\rVert_2
+
\beta_t\lVert\tilde v_t-v_t^\star\rVert_2 .
\]
Substituting $e_t=\lVert h_t-h_t^\star\rVert_2$ and the binding perturbation bound yields
\[
e_{t+1}
\le
\rho_t e_t+\beta_t(1-B_t)D_t .
\]

It remains to unroll the recurrence.
Repeated substitution gives
\[
\begin{aligned}
e_T
\le&
\left(\prod_{t=0}^{T-1}\rho_t\right)e_0\\
&+
\sum_{s=0}^{T-1}
\left(\prod_{u=s+1}^{T-1}\rho_u\right)
\beta_s(1-B_s)D_s ,
\end{aligned}
\]
which proves the claim.

\section{ATM-Bench Details}
\label{app:sec:atm-details}

This appendix will document the ATM-Bench splits, memory processors, answer models, judge models, evidence conditions, prompts, and scoring scripts used in the final experiments.

\section{STALE Details}
\label{app:sec:stale-details}

This appendix will document the STALE scenario types, probe dimensions, state-validity roles, primitive-stack implementation, and scoring protocol.

\section{Error Analysis}
\label{app:sec:error-analysis}

This appendix will contain qualitative examples grouped by substrate failure, access failure, selection failure, view-construction failure, and decoder failure.

\section{Additional Results}
\label{app:sec:additional-results}

This appendix will contain additional tables, robustness checks, judge agreement analyses, and ablations.
